\chapter{Construction by Gluing}
\label{ch:iv25}

Gluing constructs Riemann surfaces by identifying open pieces through holomorphic transition maps. This is the global counterpart of defining a manifold from compatible charts.

\section*{Learning goals}
After this chapter, the reader should be able to:
\begin{itemize}
\item construct Riemann surfaces by gluing charts along biholomorphic overlap maps;
\item verify cocycle compatibility of transition functions on multiple overlaps;
\item check Hausdorff and manifold conditions in explicit gluing constructions;
\item build standard surfaces from polygons, annuli, or slit planes;
\item track how holomorphic functions and differentials transform across glued charts;
\item distinguish equivalent gluings from genuinely different complex structures;
\end{itemize}
\section*{Conceptual roadmap}
\[
\boxed{\text{local holomorphic structure}\;\longrightarrow\;\text{integral or series representation}\;\longrightarrow\;\text{global consequence}.}
\]

\section{Gluing data}
Start with complex domains and specify identifications between open subsets.

% BEGIN VOL04-EXPANSION IV25-example-01
\begin{example}[Gluing two charts into the sphere]\label{ex:iv25-expand-01}
Take two copies \(U_0,U_\infty\cong\mathbb C\) with coordinates \(z\) and \(w\). Identify nonzero points by
\[
w=\frac1z.
\]
The transition map is biholomorphic on \(\mathbb C^*\) and is its own inverse. The quotient is the Riemann sphere: \(U_0\) covers all finite points, while \(w=0\) in the second chart represents infinity.
\end{example}
% END VOL04-EXPANSION IV25-example-01
\section{Transition maps}
Identifications must be biholomorphic and satisfy inverse and cocycle compatibility.

\section{Quotient space}
The disjoint union modulo the identifications becomes the candidate surface.

\section{Hausdorff condition}
Gluing must be controlled so the quotient remains Hausdorff.

% BEGIN VOL04-EXPANSION IV25-example-02
\begin{example}[A torus from a parallelogram]\label{ex:iv25-expand-02}
Let \(P\) be a parallelogram generated by noncollinear periods \(\omega_1,\omega_2\). Identify opposite sides by translations \(z\mapsto z+\omega_1\) and \(z\mapsto z+\omega_2\). The transition maps are holomorphic translations and the corner identifications are compatible. The quotient is topologically a torus and inherits a complex structure from the planar coordinate.
\end{example}
% END VOL04-EXPANSION IV25-example-02
\section{Charts from pieces}
Each original domain descends locally to a coordinate chart on the quotient.

\section{Sphere from two charts}
Two copies of the plane glued by $w=1/z$ on punctured regions produce the Riemann sphere.

\section{Algebraic curves}
Local branches of an algebraic relation can be glued into a smooth Riemann surface.

% BEGIN VOL04-EXPANSION IV25-example-03
\begin{example}[Two slit planes for a square root]\label{ex:iv25-expand-03}
Take two copies of the plane slit along the nonnegative real axis. Glue the upper bank of the first copy to the lower bank of the second and vice versa. Moving around the origin switches sheets, while two turns return to the starting sheet. The resulting surface carries a single-valued coordinate \(w\) with projection \(z=w^2\).
\end{example}
% END VOL04-EXPANSION IV25-example-03
\section{Cut-and-paste constructions}
Copies of slit planes or polygons can be glued to produce root surfaces and tori.

\section{Transition-cocycle viewpoint}
The global complex structure is completely encoded by holomorphic chart transition functions.

\section{Core structural results}

\begin{theorem}[Holomorphic gluing theorem]\label{thm:iv25-01}
If Hausdorff quotient gluing data have biholomorphic transition maps satisfying the cocycle conditions, the quotient inherits a unique Riemann-surface structure for which the piece maps are local charts.
\begin{proof}
Use the images of the original domains as charts. On overlaps their transition maps are exactly the prescribed biholomorphisms, so the atlas is compatible and defines the complex structure.
\end{proof}
\end{theorem}

\begin{theorem}[Sphere gluing]\label{thm:iv25-02}
Gluing two copies of $\mathbb C$ by $w=1/z$ on $\mathbb C^*$ produces a compact Riemann surface biholomorphic to $\widehat{\mathbb C}$.
\begin{proof}
Map one chart to finite complex numbers and the origin of the second chart to infinity. The transition relation is precisely the standard coordinate change at infinity.
\end{proof}
\end{theorem}

\begin{theorem}[Cocycle consistency]\label{thm:iv25-03}
For triple overlaps, transition maps must satisfy $g_{ij}\circ g_{jk}=g_{ik}$.
\begin{proof}
A point represented through three charts must have the same coordinate identification independent of the path through the overlap graph; the cocycle condition is exactly that independence.
\end{proof}
\end{theorem}

\section{Worked examples}

\begin{example}[Gluing data diagnostic]\label{ex:iv25-worked-01}
Give a rigorous worked analysis of Gluing data. State the relevant holomorphic, contour, series, or singularity hypothesis and derive the central conclusion. Start with complex domains and specify identifications between open subsets. The decisive step is to use the complex-analytic structure globally rather than reason from real-variable intuition alone.
\end{example}

\begin{example}[Transition maps diagnostic]\label{ex:iv25-worked-02}
Give a rigorous worked analysis of Transition maps. State the relevant holomorphic, contour, series, or singularity hypothesis and derive the central conclusion. Identifications must be biholomorphic and satisfy inverse and cocycle compatibility. The decisive step is to use the complex-analytic structure globally rather than reason from real-variable intuition alone.
\end{example}

\begin{example}[Quotient space diagnostic]\label{ex:iv25-worked-03}
Give a rigorous worked analysis of Quotient space. State the relevant holomorphic, contour, series, or singularity hypothesis and derive the central conclusion. The disjoint union modulo the identifications becomes the candidate surface. The decisive step is to use the complex-analytic structure globally rather than reason from real-variable intuition alone.
\end{example}

\begin{example}[Hausdorff condition diagnostic]\label{ex:iv25-worked-04}
Give a rigorous worked analysis of Hausdorff condition. State the relevant holomorphic, contour, series, or singularity hypothesis and derive the central conclusion. Gluing must be controlled so the quotient remains Hausdorff. The decisive step is to use the complex-analytic structure globally rather than reason from real-variable intuition alone.
\end{example}

\section{Solved dossiers}
The dossiers are canonical solved problems. The provenance ledger separately records which retained corpus rules guided each topic and which dossiers were newly designed.

\begin{problem}[Gluing data diagnostic]\label{prob:iv25-dossier-01}
Give a rigorous worked analysis of Gluing data. State the relevant holomorphic, contour, series, or singularity hypothesis and derive the central conclusion.
\end{problem}
\begin{solution}
Start with complex domains and specify identifications between open subsets. The decisive step is to use the complex-analytic structure globally rather than reason from real-variable intuition alone.
\end{solution}

\begin{problem}[Transition maps diagnostic]\label{prob:iv25-dossier-02}
Give a rigorous worked analysis of Transition maps. State the relevant holomorphic, contour, series, or singularity hypothesis and derive the central conclusion.
\end{problem}
\begin{solution}
Identifications must be biholomorphic and satisfy inverse and cocycle compatibility. The decisive step is to use the complex-analytic structure globally rather than reason from real-variable intuition alone.
\end{solution}

\begin{problem}[Quotient space diagnostic]\label{prob:iv25-dossier-03}
Give a rigorous worked analysis of Quotient space. State the relevant holomorphic, contour, series, or singularity hypothesis and derive the central conclusion.
\end{problem}
\begin{solution}
The disjoint union modulo the identifications becomes the candidate surface. The decisive step is to use the complex-analytic structure globally rather than reason from real-variable intuition alone.
\end{solution}

\begin{problem}[Hausdorff condition diagnostic]\label{prob:iv25-dossier-04}
Give a rigorous worked analysis of Hausdorff condition. State the relevant holomorphic, contour, series, or singularity hypothesis and derive the central conclusion.
\end{problem}
\begin{solution}
Gluing must be controlled so the quotient remains Hausdorff. The decisive step is to use the complex-analytic structure globally rather than reason from real-variable intuition alone.
\end{solution}

\begin{problem}[Charts from pieces diagnostic]\label{prob:iv25-dossier-05}
Give a rigorous worked analysis of Charts from pieces. State the relevant holomorphic, contour, series, or singularity hypothesis and derive the central conclusion.
\end{problem}
\begin{solution}
Each original domain descends locally to a coordinate chart on the quotient. The decisive step is to use the complex-analytic structure globally rather than reason from real-variable intuition alone.
\end{solution}

\begin{problem}[Sphere from two charts diagnostic]\label{prob:iv25-dossier-06}
Give a rigorous worked analysis of Sphere from two charts. State the relevant holomorphic, contour, series, or singularity hypothesis and derive the central conclusion.
\end{problem}
\begin{solution}
Two copies of the plane glued by $w=1/z$ on punctured regions produce the Riemann sphere. The decisive step is to use the complex-analytic structure globally rather than reason from real-variable intuition alone.
\end{solution}

\begin{problem}[Algebraic curves diagnostic]\label{prob:iv25-dossier-07}
Give a rigorous worked analysis of Algebraic curves. State the relevant holomorphic, contour, series, or singularity hypothesis and derive the central conclusion.
\end{problem}
\begin{solution}
Local branches of an algebraic relation can be glued into a smooth Riemann surface. The decisive step is to use the complex-analytic structure globally rather than reason from real-variable intuition alone.
\end{solution}

\begin{problem}[Cut-and-paste constructions diagnostic]\label{prob:iv25-dossier-08}
Give a rigorous worked analysis of Cut-and-paste constructions. State the relevant holomorphic, contour, series, or singularity hypothesis and derive the central conclusion.
\end{problem}
\begin{solution}
Copies of slit planes or polygons can be glued to produce root surfaces and tori. The decisive step is to use the complex-analytic structure globally rather than reason from real-variable intuition alone.
\end{solution}

\begin{problem}[Holomorphic gluing theorem proof dossier]\label{prob:iv25-dossier-09}
Prove the following theorem-level statement and identify the essential hypothesis: If Hausdorff quotient gluing data have biholomorphic transition maps satisfying the cocycle conditions, the quotient inherits a unique Riemann-surface structure for which the piece maps are local charts.
\end{problem}
\begin{solution}
Use the images of the original domains as charts. On overlaps their transition maps are exactly the prescribed biholomorphisms, so the atlas is compatible and defines the complex structure. This identifies the mechanism that makes the complex-analytic conclusion stronger than its real-variable analogue.
\end{solution}

\begin{problem}[Sphere gluing proof dossier]\label{prob:iv25-dossier-10}
Prove the following theorem-level statement and identify the essential hypothesis: Gluing two copies of $\mathbb C$ by $w=1/z$ on $\mathbb C^*$ produces a compact Riemann surface biholomorphic to $\widehat{\mathbb C}$.
\end{problem}
\begin{solution}
Map one chart to finite complex numbers and the origin of the second chart to infinity. The transition relation is precisely the standard coordinate change at infinity. This identifies the mechanism that makes the complex-analytic conclusion stronger than its real-variable analogue.
\end{solution}

\begin{problem}[Cocycle consistency proof dossier]\label{prob:iv25-dossier-11}
Prove the following theorem-level statement and identify the essential hypothesis: For triple overlaps, transition maps must satisfy $g_{ij}\circ g_{jk}=g_{ik}$.
\end{problem}
\begin{solution}
A point represented through three charts must have the same coordinate identification independent of the path through the overlap graph; the cocycle condition is exactly that independence. This identifies the mechanism that makes the complex-analytic conclusion stronger than its real-variable analogue.
\end{solution}

\begin{problem}[Chapter synthesis]\label{prob:iv25-dossier-12}
Connect the definitions and principal theorems of this chapter into one reusable complex-analysis strategy.
\end{problem}
\begin{solution}
Gluing constructs Riemann surfaces by identifying open pieces through holomorphic transition maps. This is the global counterpart of defining a manifold from compatible charts. The reusable strategy is to identify the analytic domain, choose the correct local expansion or contour representation, apply the structural theorem, and then audit singularities and topology.
\end{solution}

\section{Exercises with complete solutions}

\begin{exercise}\label{exr:iv25-01}
State and apply the central principle behind Gluing data. Give one concrete consequence.
\end{exercise}
\begin{hint}
Write every overlap map explicitly and check it is biholomorphic on the overlap region.
\end{hint}
\begin{solution}
Start with complex domains and specify identifications between open subsets. The same principle can then be reused in later contour, residue, conformal, or Riemann-surface arguments.
\end{solution}

\begin{exercise}\label{exr:iv25-02}
State and apply the central principle behind Transition maps. Give one concrete consequence.
\end{exercise}
\begin{hint}
On triple overlaps, verify \(g_{ij}g_{jk}=g_{ik}\); this cocycle identity makes chart identifications consistent.
\end{hint}
\begin{solution}
Identifications must be biholomorphic and satisfy inverse and cocycle compatibility. The same principle can then be reused in later contour, residue, conformal, or Riemann-surface arguments.
\end{solution}

\begin{exercise}\label{exr:iv25-03}
State and apply the central principle behind Quotient space. Give one concrete consequence.
\end{exercise}
\begin{hint}
To check Hausdorffness, separate inequivalent points using neighborhoods that survive the quotient.
\end{hint}
\begin{solution}
The disjoint union modulo the identifications becomes the candidate surface. The same principle can then be reused in later contour, residue, conformal, or Riemann-surface arguments.
\end{solution}

\begin{exercise}\label{exr:iv25-04}
State and apply the central principle behind Hausdorff condition. Give one concrete consequence.
\end{exercise}
\begin{hint}
Polygon edge identifications should be orientation-compatible with the intended complex structure; track the local coordinates near vertices.
\end{hint}
\begin{solution}
Gluing must be controlled so the quotient remains Hausdorff. The same principle can then be reused in later contour, residue, conformal, or Riemann-surface arguments.
\end{solution}

\begin{exercise}\label{exr:iv25-05}
State and apply the central principle behind Charts from pieces. Give one concrete consequence.
\end{exercise}
\begin{hint}
Functions glue only when their local expressions agree after composing with transition maps.
\end{hint}
\begin{solution}
Each original domain descends locally to a coordinate chart on the quotient. The same principle can then be reused in later contour, residue, conformal, or Riemann-surface arguments.
\end{solution}

\begin{exercise}\label{exr:iv25-06}
State and apply the central principle behind Sphere from two charts. Give one concrete consequence.
\end{exercise}
\begin{hint}
Differentials glue with the derivative of the transition function; check the transformation law on overlaps.
\end{hint}
\begin{solution}
Two copies of the plane glued by $w=1/z$ on punctured regions produce the Riemann sphere. The same principle can then be reused in later contour, residue, conformal, or Riemann-surface arguments.
\end{solution}

\begin{exercise}\label{exr:iv25-07}
State and apply the central principle behind Algebraic curves. Give one concrete consequence.
\end{exercise}
\begin{hint}
Equivalent refinements or chart changes do not alter the underlying Riemann surface; compare transition atlases.
\end{hint}
\begin{solution}
Local branches of an algebraic relation can be glued into a smooth Riemann surface. The same principle can then be reused in later contour, residue, conformal, or Riemann-surface arguments.
\end{solution}

\begin{exercise}\label{exr:iv25-08}
State and apply the central principle behind Cut-and-paste constructions. Give one concrete consequence.
\end{exercise}
\begin{hint}
Watch for quotient identifications that create nonmanifold points; local neighborhoods must still look like disks.
\end{hint}
\begin{solution}
Copies of slit planes or polygons can be glued to produce root surfaces and tori. The same principle can then be reused in later contour, residue, conformal, or Riemann-surface arguments.
\end{solution}

\section*{Chapter summary}
The chapter now forms part of the canonical complex-analysis chain from holomorphicity through residues, global continuation, and Riemann surfaces.
% BEGIN VOL04-EXPANSION IV25-exercises-01
\section{Graded supplementary exercises}
These exercises extend the protected chapter with a balanced set of computations, proofs, hypothesis tests, applications, and challenges.

\subsection*{Standard computations}

\begin{exercise}[Sphere transition]\label{exr:iv25-expand-01}
What is the transition map between the finite and infinite charts of the sphere?
\end{exercise}
\begin{hint}
Use the reciprocal coordinate.
\end{hint}
\begin{solution}
It is \(w=1/z\) on \(\mathbb C^*\).
\end{solution}

\begin{exercise}[Inverse transition]\label{exr:iv25-expand-02}
What is the inverse of \(w=1/z\)?
\end{exercise}
\begin{hint}
The reciprocal map is an involution.
\end{hint}
\begin{solution}
It is \(z=1/w\).
\end{solution}

\begin{exercise}[Torus edge map]\label{exr:iv25-expand-03}
What holomorphic map identifies opposite vertical edges of a lattice parallelogram?
\end{exercise}
\begin{hint}
Use translation by one period.
\end{hint}
\begin{solution}
It is a translation \(z\mapsto z+\omega_1\) or its inverse, depending on the chosen edge orientation.
\end{solution}

\begin{exercise}[Cocycle check]\label{exr:iv25-expand-04}
If \(g_{12}(z)=z+1\) and \(g_{23}(z)=z+i\), what must \(g_{13}\) be on a triple overlap?
\end{exercise}
\begin{hint}
Compose the first two transitions.
\end{hint}
\begin{solution}
The cocycle condition gives \(g_{13}(z)=z+1+i\).
\end{solution}

\begin{exercise}[Square-root sheet switch]\label{exr:iv25-expand-05}
How many crossings of the slit are needed to return to the original sheet in the two-sheeted square-root gluing?
\end{exercise}
\begin{hint}
Each crossing switches sheets.
\end{hint}
\begin{solution}
Two crossings return to the original sheet.
\end{solution}

\subsection*{Proofs}

\begin{exercise}[Gluing theorem]\label{exr:iv25-expand-06}
Explain why compatible biholomorphic transition maps define a complex structure on a Hausdorff quotient.
\end{exercise}
\begin{hint}
Use the images of the original pieces as charts.
\end{hint}
\begin{solution}
Each piece descends locally homeomorphically to the quotient, and on overlaps the chart transitions are exactly the prescribed biholomorphisms. Thus the atlas is holomorphically compatible.
\end{solution}

\begin{exercise}[Sphere compactness]\label{exr:iv25-expand-07}
Why is the two-chart sphere gluing compact?
\end{exercise}
\begin{hint}
Relate it to the one-point compactification of the plane.
\end{hint}
\begin{solution}
The second chart provides a neighborhood of the added infinity point, and the resulting quotient is homeomorphic to the compact Riemann sphere.
\end{solution}

\begin{exercise}[Cocycle necessity]\label{exr:iv25-expand-08}
Why is the cocycle condition necessary on triple overlaps?
\end{exercise}
\begin{hint}
A point can be transferred between charts by two different routes.
\end{hint}
\begin{solution}
Without \(g_{ij}\circ g_{jk}=g_{ik}\), the two routes assign different coordinates to the same equivalence class, so the quotient atlas is inconsistent.
\end{solution}

\begin{exercise}[Translation transitions]\label{exr:iv25-expand-09}
Prove that a lattice quotient inherits a Riemann-surface structure.
\end{exercise}
\begin{hint}
Use disks smaller than the shortest nonzero lattice displacement.
\end{hint}
\begin{solution}
Such a disk has disjoint lattice translates. Projection is one to one on it, and any overlap transition between lifted disks is a translation, hence holomorphic.
\end{solution}

\subsection*{Counterexamples and hypothesis tests}

\begin{exercise}[Noninvertible transition]\label{exr:iv25-expand-10}
Can \(w=z^2\) be used as a chart transition across a neighborhood containing zero?
\end{exercise}
\begin{hint}
A transition must be biholomorphic.
\end{hint}
\begin{solution}
No. Near zero the map is not locally one to one and its derivative vanishes.
\end{solution}

\begin{exercise}[Non-Hausdorff quotient]\label{exr:iv25-expand-11}
Is holomorphic compatibility alone enough if the quotient topology is non-Hausdorff?
\end{exercise}
\begin{hint}
Recall the manifold separation axiom.
\end{hint}
\begin{solution}
No. A Riemann surface must be Hausdorff, so pathological gluing can fail before complex analysis is considered.
\end{solution}

\begin{exercise}[Arbitrary boundary identification]\label{exr:iv25-expand-12}
May opposite polygon edges be glued by any continuous map and still produce a Riemann surface?
\end{exercise}
\begin{hint}
The transition must respect complex structure.
\end{hint}
\begin{solution}
No. The local identifications used as chart transitions must be holomorphic or antiholomorphic only if one is deliberately changing orientation; for a Riemann surface atlas they must be biholomorphic.
\end{solution}

\subsection*{Applications and investigations}

\begin{exercise}[Polygon models]\label{exr:iv25-expand-13}
Why are polygon gluings useful for constructing compact surfaces?
\end{exercise}
\begin{hint}
They encode global topology through finitely many edge identifications.
\end{hint}
\begin{solution}
A single polygon with paired edges can represent the whole surface, while affine or translational edge maps often make the induced complex charts explicit.
\end{solution}

\begin{exercise}[Algebraic curve normalization]\label{exr:iv25-expand-14}
How can gluing local branches help construct a smooth surface from an algebraic relation?
\end{exercise}
\begin{hint}
Use separate local parameters near branch behavior.
\end{hint}
\begin{solution}
Smooth local branches are treated as charts and glued where their coordinate descriptions agree, separating sheets that appear multivalued in the plane projection.
\end{solution}

\subsection*{Challenge problems}

\begin{exercise}[Sphere from disks]\label{exr:iv25-expand-15}
Construct the sphere by gluing two unit disks along annular collars rather than two full planes.
\end{exercise}
\begin{hint}
Use reciprocal coordinates on the overlap annuli.
\end{hint}
\begin{solution}
Choose annuli where \(w=1/z\) maps one collar biholomorphically to the other. The resulting compact quotient has the same two-chart complex structure as the sphere.
\end{solution}

\begin{exercise}[Corner consistency on a torus]\label{exr:iv25-expand-16}
Why do the four vertices of a fundamental parallelogram become one point without creating a singularity?
\end{exercise}
\begin{hint}
Use translations around the corner.
\end{hint}
\begin{solution}
Neighborhood sectors from the four corners glue by translations into one ordinary disk neighborhood. The total angle and holomorphic coordinate match smoothly, so no cone singularity appears.
\end{solution}

% END VOL04-EXPANSION IV25-exercises-01
